\documentclass[11pt]{article}
\usepackage[margin=1in]{geometry}
\usepackage[T1]{fontenc}
\usepackage{lmodern}
\usepackage{amsmath,amssymb}
\setlength{\parindent}{0pt}
\setlength{\parskip}{0.8em}
\title{Linear Algebra: Abstract Structures, Eigentheory, and Spectral Analysis}
\author{Advanced Course Notes}
\date{}
\begin{document}
\maketitle

\section*{Abstract Vector Spaces}

A vector space over a field $F$ is a set $V$ together with two operations, addition $+: V \times V \to V$ and scalar multiplication $\cdot: F \times V \to V$, satisfying eight axioms. For $u, v, w \in V$ and $a, b \in F$:
\begin{enumerate}
  \item $u + v = v + u$ (commutativity).
  \item $(u + v) + w = u + (v + w)$ (associativity).
  \item There exists $\mathbf{0} \in V$ such that $v + \mathbf{0} = v$ for all $v$.
  \item For each $v$ there exists $-v$ such that $v + (-v) = \mathbf{0}$.
  \item $1 \cdot v = v$ (identity scalar).
  \item $a(bv) = (ab)v$ (associativity of scalar multiplication).
  \item $a(u + v) = au + av$ (distributivity over vector addition).
  \item $(a + b)v = av + bv$ (distributivity over scalar addition).
\end{enumerate}

The key insight is that $V$ need not consist of tuples of numbers. Polynomials of degree at most $n$, continuous functions on an interval, infinite sequences, and matrices of fixed size are all vector spaces under appropriate operations. The axioms capture the structural essence of ``linear combination'' divorced from any particular representation.

A subspace $W \leq V$ is a subset closed under addition and scalar multiplication (equivalently, closed under all linear combinations). The empty set fails to be a subspace since it lacks $\mathbf{0}$; the canonical subspaces are $\{\mathbf{0}\}$ and $V$ itself.

The intersection of any family of subspaces is a subspace. The sum $W_1 + W_2 = \{w_1 + w_2 : w_1 \in W_1, w_2 \in W_2\}$ is a subspace. If $W_1 \cap W_2 = \{\mathbf{0}\}$, the sum is direct: $V = W_1 \oplus W_2$, and every $v \in V$ decomposes uniquely as $w_1 + w_2$.

\section*{Linear Independence, Span, and Dimension}

A set $S \subseteq V$ is linearly independent if no finite linear combination $\sum_i c_i v_i = \mathbf{0}$ (with $v_i \in S$) is possible except with all coefficients zero. This is equivalent to: no element of $S$ lies in the span of the others.

The span of $S$ is the smallest subspace containing $S$: all finite linear combinations of elements of $S$. A set $S$ that is both independent and spanning for $V$ is a basis. By the replacement theorem (or Steinitz exchange lemma), any two bases of $V$ have the same cardinality, called the dimension $\dim V$.

In a finite-dimensional space of dimension $n$:
\begin{itemize}
  \item Any independent set has at most $n$ elements; if it has exactly $n$, it is a basis.
  \item Any spanning set has at least $n$ elements; if it has exactly $n$, it is a basis.
  \item Any $n$ independent vectors automatically span $V$.
\end{itemize}

The coordinate isomorphism: once a basis $\mathcal{B} = (b_1, \ldots, b_n)$ is fixed, every $v = \sum_i c_i b_i$ is represented by its coordinate vector $[v]_\mathcal{B} = (c_1, \ldots, c_n)^T \in \mathbb{R}^n$. This isomorphism is canonical given $\mathcal{B}$ but depends on the ordering.

\section*{Linear Maps: Kernel, Image, Rank-Nullity}

A linear map $T: V \to W$ between vector spaces satisfies $T(au + bv) = aT(u) + bT(v)$ for all $u,v \in V$ and $a,b \in F$. Linear maps are morphisms in the category of vector spaces: they preserve the algebraic structure.

The kernel $\ker T = \{v \in V : T(v) = \mathbf{0}\}$ is a subspace of $V$. The image $\text{im}\, T = \{T(v) : v \in V\}$ is a subspace of $W$.

$T$ is injective if and only if $\ker T = \{\mathbf{0}\}$. $T$ is surjective if and only if $\text{im}\, T = W$. A bijective linear map is an isomorphism.

The rank-nullity theorem: for $T: V \to W$ with $V$ finite-dimensional,
\[
\dim \ker T + \dim \text{im}\, T = \dim V.
\]
This fundamental identity constrains how a linear map can distribute the dimension of its domain between what it kills and what it produces. It implies that $T$ cannot be injective and fail to be surjective simultaneously when $\dim V = \dim W$.

The matrix of $T$ with respect to bases $\mathcal{B}$ (for $V$) and $\mathcal{C}$ (for $W$) is the $m \times n$ matrix $[T]_{\mathcal{B}}^{\mathcal{C}}$ whose $j$-th column is $[T(b_j)]_\mathcal{C}$. Changing bases replaces this matrix by $Q^{-1} [T]_{\mathcal{B}}^{\mathcal{C}} P$, where $P$ and $Q$ are the change-of-basis matrices.

\section*{Eigenvalues and Eigenvectors: Abstract Formulation}

Let $T: V \to V$ be a linear operator on a finite-dimensional vector space $V$ over $F$. A scalar $\lambda \in F$ is an eigenvalue of $T$ if there exists a nonzero $v \in V$ with
\[
T(v) = \lambda v.
\]
The nonzero vectors satisfying this are the eigenvectors corresponding to $\lambda$. The eigenspace is
\[
E_\lambda = \ker(T - \lambda I) = \{v \in V : T(v) = \lambda v\},
\]
which includes $\mathbf{0}$ and is a nonzero subspace of $V$.

The eigenvalue equation $T(v) = \lambda v$ can be rewritten as $(T - \lambda I)(v) = \mathbf{0}$. An eigenvalue exists precisely when $T - \lambda I$ is not injective, i.e., when $\ker(T - \lambda I) \neq \{\mathbf{0}\}$.

In coordinates, if $A = [T]_\mathcal{B}$ is the matrix of $T$ in some basis, then $\lambda$ is an eigenvalue of $T$ if and only if
\[
\det(A - \lambda I) = 0.
\]
The eigenvalues are coordinate-free invariants of $T$: they do not depend on the choice of basis, since a change of basis replaces $A$ by $P^{-1}AP$, and $\det(P^{-1}AP - \lambda I) = \det(P^{-1}(A-\lambda I)P) = \det(A - \lambda I)$.

\section*{The Characteristic Polynomial}

The characteristic polynomial of $T$ (or of its matrix $A$) is
\[
p_A(\lambda) = \det(A - \lambda I) = (-1)^n \lambda^n + \cdots + \det(A).
\]
This is a degree-$n$ polynomial. Its leading coefficient is $(-1)^n$ (for the $\det(-\lambda I)$ term). Two notable coefficients:
\begin{itemize}
  \item Constant term ($\lambda = 0$): $p_A(0) = \det(A)$.
  \item Coefficient of $\lambda^{n-1}$: up to sign, equals $\text{tr}(A)$, the trace (sum of diagonal entries).
\end{itemize}

By the factor theorem, the roots of $p_A(\lambda)$ are the eigenvalues. Over $\mathbb{C}$, the fundamental theorem of algebra guarantees $n$ roots (counted with multiplicity). Over $\mathbb{R}$, some eigenvalues may be complex.

The Cayley-Hamilton theorem states that every matrix satisfies its own characteristic polynomial: $p_A(A) = 0$. This is a remarkable result: substituting the matrix itself into the polynomial gives the zero matrix. It has practical implications, including that $A^{-1}$ can be expressed as a polynomial in $A$ (of degree at most $n-1$) when $A$ is invertible.

The algebraic multiplicity $m_a(\lambda)$ of an eigenvalue is its multiplicity as a root of $p_A$. The geometric multiplicity $m_g(\lambda) = \dim E_\lambda = \dim \ker(A - \lambda I)$ satisfies $1 \leq m_g(\lambda) \leq m_a(\lambda)$.

\section*{Diagonalization and the Spectral Decomposition}

A linear operator $T: V \to V$ is diagonalizable if there exists a basis of $V$ consisting entirely of eigenvectors. In this eigenbasis, the matrix of $T$ is diagonal:
\[
[T]_\mathcal{E} = \text{diag}(\lambda_1, \lambda_2, \ldots, \lambda_n).
\]
In terms of matrices: $A$ is diagonalizable if there exists an invertible $P$ with $P^{-1}AP = D$ diagonal.

Diagonalizability criteria:
\begin{itemize}
  \item $A$ is diagonalizable iff $m_g(\lambda) = m_a(\lambda)$ for every eigenvalue $\lambda$.
  \item Equivalently, the eigenspaces together have total dimension $n$.
  \item If $A$ has $n$ distinct eigenvalues, it is automatically diagonalizable (since eigenvectors for distinct eigenvalues are independent).
\end{itemize}

Eigenvectors corresponding to distinct eigenvalues are linearly independent. Proof: suppose $T(v_i) = \lambda_i v_i$ for distinct $\lambda_i$ and $\sum c_i v_i = 0$. Apply $T - \lambda_1 I$ to eliminate $v_1$, giving $\sum_{i>1} c_i (\lambda_i - \lambda_1) v_i = 0$. Repeating eliminates vectors one by one; since the $\lambda_i$ are distinct and the $v_i$ are nonzero, all $c_i = 0$.

The spectral decomposition: if $A = PDP^{-1}$ with $D = \text{diag}(\lambda_1, \ldots, \lambda_n)$, writing $P = [v_1 \mid \cdots \mid v_n]$ and $P^{-1} = \begin{bmatrix} w_1^T \\ \vdots \\ w_n^T \end{bmatrix}$:
\[
A = \sum_{i=1}^n \lambda_i v_i w_i^T.
\]
Each $v_i w_i^T$ is a rank-one projector onto the eigenspace of $\lambda_i$ along the complementary space.

\section*{Vector Spaces with Inner Products}

An inner product on a real vector space $V$ is a map $\langle \cdot, \cdot \rangle : V \times V \to \mathbb{R}$ satisfying:
\begin{enumerate}
  \item Symmetry: $\langle u, v \rangle = \langle v, u \rangle$.
  \item Linearity in first argument: $\langle au + bv, w \rangle = a\langle u, w\rangle + b\langle v, w\rangle$.
  \item Positive definiteness: $\langle v, v \rangle \geq 0$ with equality iff $v = \mathbf{0}$.
\end{enumerate}

The induced norm is $\|v\| = \sqrt{\langle v, v \rangle}$. The Cauchy-Schwarz inequality states $|\langle u, v \rangle| \leq \|u\| \|v\|$, with equality iff $u$ and $v$ are proportional. The triangle inequality follows: $\|u + v\| \leq \|u\| + \|v\|$.

Orthogonality in the inner product sense: $u \perp v$ iff $\langle u, v \rangle = 0$. An orthonormal basis $\{q_1, \ldots, q_n\}$ satisfies $\langle q_i, q_j \rangle = \delta_{ij}$. In orthonormal coordinates, inner products and matrix representations simplify considerably.

The Gram-Schmidt process: given a basis $\{v_1, \ldots, v_n\}$, produce orthonormal $\{q_1, \ldots, q_n\}$ by
\[
u_k = v_k - \sum_{j=1}^{k-1} \langle v_k, q_j \rangle q_j, \qquad q_k = \frac{u_k}{\|u_k\|}.
\]
This is the orthogonal projection onto the orthogonal complement of the span of the previous vectors, followed by normalization.

\section*{Orthogonality: Subspaces and Projections}

The orthogonal complement of a subspace $W \leq V$ is
\[
W^\perp = \{v \in V : \langle v, w \rangle = 0 \text{ for all } w \in W\}.
\]
This is a subspace, $\dim W + \dim W^\perp = \dim V$, $(W^\perp)^\perp = W$, and $V = W \oplus W^\perp$ (orthogonal direct sum).

The orthogonal projection $P_W: V \to W$ sends each $v$ to its nearest point in $W$ (in the inner product norm). It is characterized by $P_W v \in W$ and $v - P_W v \in W^\perp$. In terms of an orthonormal basis $\{q_1, \ldots, q_k\}$ for $W$:
\[
P_W v = \sum_{i=1}^k \langle v, q_i \rangle q_i.
\]
As a matrix: $P = QQ^T$ where $Q = [q_1 \mid \cdots \mid q_k]$. Projections are idempotent ($P^2 = P$) and self-adjoint ($P^T = P$).

The best approximation theorem: $P_W v$ is the closest vector in $W$ to $v$, i.e., $\|v - P_W v\| \leq \|v - w\|$ for all $w \in W$. This underpins least-squares regression, Fourier analysis, and signal decomposition.

\section*{Symmetric Matrices and the Spectral Theorem}

A matrix $A$ is symmetric if $A = A^T$ (equivalently, $A_{ij} = A_{ji}$). Symmetric matrices are the real analogues of Hermitian operators and have especially clean eigentheory.

The real spectral theorem: if $A$ is a real symmetric $n \times n$ matrix, then:
\begin{enumerate}
  \item All eigenvalues of $A$ are real.
  \item Eigenvectors corresponding to distinct eigenvalues are orthogonal.
  \item $A$ is orthogonally diagonalizable: there exists an orthogonal matrix $Q$ (with $Q^{-1} = Q^T$) such that $Q^T A Q = D$ is diagonal.
\end{enumerate}

Proof of orthogonality of eigenvectors: suppose $Av = \lambda v$ and $Aw = \mu w$ with $\lambda \neq \mu$. Then
\[
\lambda \langle v, w \rangle = \langle Av, w \rangle = \langle v, A^T w \rangle = \langle v, Aw \rangle = \mu \langle v, w \rangle.
\]
Since $\lambda \neq \mu$, we conclude $\langle v, w \rangle = 0$.

The spectral decomposition for symmetric matrices: $A = \sum_{i=1}^n \lambda_i q_i q_i^T$, where $\{q_i\}$ is an orthonormal eigenbasis. Each $q_i q_i^T$ is the orthogonal projection onto the span of $q_i$.

\section*{Eigenvalues and Determinants: The Structural Connection}

The determinant $\det(A)$ equals the product of all eigenvalues (over $\mathbb{C}$, counted with algebraic multiplicity):
\[
\det(A) = \lambda_1 \lambda_2 \cdots \lambda_n.
\]
This follows because $\det(A) = p_A(0) = \prod_i (0 - \lambda_i) \cdot (-1)^n = \prod_i \lambda_i$ (adjusting for the sign convention).

The trace $\text{tr}(A) = \sum_{i=1}^n A_{ii}$ equals the sum of eigenvalues:
\[
\text{tr}(A) = \lambda_1 + \lambda_2 + \cdots + \lambda_n.
\]
Both trace and determinant are invariants under similarity: if $B = P^{-1}AP$, then $\text{tr}(B) = \text{tr}(A)$ and $\det(B) = \det(A)$. The entire characteristic polynomial is a similarity invariant.

A matrix is invertible iff 0 is not an eigenvalue (since $\det(A) = 0$ iff some $\lambda_i = 0$). The inverse has eigenvalues $\{1/\lambda_i\}$ with the same eigenvectors.

For a matrix $A$ with eigenvalue $\lambda$ and eigenvector $v$: $A^k v = \lambda^k v$, and if $f$ is any polynomial, $f(A) v = f(\lambda) v$. This is the key computational leverage of eigenvectors: they diagonalize polynomial functions of the matrix.

\section*{Generalized Eigenvectors and Jordan Normal Form}

When $A$ is not diagonalizable (some $m_g(\lambda) < m_a(\lambda)$), there is still a canonical form: the Jordan normal form. A generalized eigenvector of rank $k$ satisfies
\[
(A - \lambda I)^k v = \mathbf{0}, \quad (A - \lambda I)^{k-1} v \neq \mathbf{0}.
\]

For each eigenvalue $\lambda$, the generalized eigenspace $\ker(A - \lambda I)^{m_a(\lambda)}$ has dimension $m_a(\lambda)$. Together the generalized eigenspaces decompose $\mathbb{C}^n$ into a direct sum (the primary decomposition).

The Jordan normal form $J$ of $A$ is block diagonal, with one Jordan block per ``chain'' of generalized eigenvectors:
\[
J_k(\lambda) = \begin{bmatrix} \lambda & 1 & & \\ & \lambda & \ddots & \\ & & \ddots & 1 \\ & & & \lambda \end{bmatrix} \in \mathbb{R}^{k \times k}.
\]
Diagonal matrices are Jordan normal forms where every block has size 1.

\section*{Matrix Inverse and Its Relationship to Eigenvalues}

If $A$ is invertible, then $A^{-1}$ has eigenvalues $1/\lambda_i$ with the same eigenvectors as $A$. This follows directly: $Av = \lambda v \Rightarrow v = \lambda A^{-1} v \Rightarrow A^{-1} v = (1/\lambda) v$.

The Cayley-Hamilton theorem provides a formula for $A^{-1}$: since $p_A(A) = 0$ and $p_A(\lambda) = (-1)^n \lambda^n + \cdots + \det(A)$, dividing by $\det(A)$ and solving for $I$ gives
\[
A^{-1} = \frac{-1}{\det(A)} \bigl[(-1)^n A^{n-1} + \cdots + c_1 I\bigr],
\]
expressing $A^{-1}$ as a polynomial in $A$. This is theoretically important but seldom used for numerical computation.

For the matrix exponential $e^A = \sum_{k=0}^\infty A^k / k!$, diagonalization gives $e^A = P e^D P^{-1}$ where $e^D = \text{diag}(e^{\lambda_1}, \ldots, e^{\lambda_n})$. The matrix exponential arises in systems of differential equations: the solution to $\dot{x} = Ax$ with $x(0) = x_0$ is $x(t) = e^{tA} x_0$.

\section*{Rank, Linear Independence, and the Fundamental Theorem}

The fundamental theorem of linear algebra describes four fundamental subspaces of an $m \times n$ matrix $A$ with rank $r$:
\begin{enumerate}
  \item Column space $C(A) \leq \mathbb{R}^m$, dimension $r$.
  \item Null space $N(A) \leq \mathbb{R}^n$, dimension $n - r$.
  \item Row space $C(A^T) \leq \mathbb{R}^n$, dimension $r$.
  \item Left null space $N(A^T) \leq \mathbb{R}^m$, dimension $m - r$.
\end{enumerate}

The orthogonality relations: $N(A) = C(A^T)^\perp$ and $N(A^T) = C(A)^\perp$. Thus $\mathbb{R}^n = C(A^T) \oplus N(A)$ and $\mathbb{R}^m = C(A) \oplus N(A^T)$ as orthogonal direct sums.

This means: any vector in the input space $\mathbb{R}^n$ decomposes uniquely into a component in the row space (which gets mapped injectively into the column space) and a component in the null space (which gets annihilated). The transformation $A$ is invertible on the row space, with image equal to the column space.

This structural picture explains why rank is the most fundamental invariant of a linear map: it simultaneously measures the dimension of what is preserved and constrains the dimension of what is lost.

\end{document}
